% a06_BALtorus_d4: part of Step 3 of the proof of thm:a06:T1 (input from sec/a06_BALtorus_c.tex).
% Group (C), first half: the large-field chain of Ba2 Sect. 3A, (3.1)-(3.20), and the k-uniformity of Ba2's unqualified O(1) lines.

\subsubsection*{Step 3(C$_1$): the large-field chain of \cite[Sect.~3A]{Ba2}}

\cite[Sect.~3A]{Ba2} inverts the renormalisation transformations in the large-field region: the inductive hypothesis \cite[(3.5)]{Ba2} at level $k$ is reduced to the same statement at level $k-1$ at the cost of a factor $\exp(O(1)(L^k\eps)^{\kappa_0}|T_\eps|)$, through the Gaussian integral \cite[(3.9)]{Ba2} with the perturbing kernels $H'_k,H''_k$ \cite[(3.8), (3.12)]{Ba2}, the linear terms $F'_k,F''_k$ \cite[(3.10)]{Ba2}, the lower bounds \cite[(3.11)]{Ba2} and the Gaussian comparison \cite[(3.17)--(3.20)]{Ba2}. What is consumed is that each inverse step costs $\exp(C_\kappa t_k^\kappa|T^{(k)}_1|)$ with $C_\kappa$ independent of $k$, $\eps$, $\eps_0$ and the torus. We prove this with explicit constants. Besides the notation of Step~3(B), $M_v\defeq c_Ap(t)/(\lmu t)$ and $M_s\defeq c_sp(t)(\lambda t)^{-1/4}$ are the small-field bounds \cite[(2.16)--(2.17)]{Ba2} at scale $k$ ($c_A\defeq2+2Ld\lmu$, $c_s\defeq2+2Ld$, under the cap $\lambda t\le1$), $M(t)\defeq c_Mt^{-1}p(t)\ge\max(M_v,M_s)$ with $c_M\defeq\max(c_A/\lmu,c_s\lambda^{-1/4})$; scale-$(k-1)$ collar fields rescaled to scale $k$ are bounded by $c_LM_v(t)$ with $c_L\defeq L^{3/2}(1+\ln L)^{\bar p}$; $S_\delta\defeq\sup_x\sum_ze^{-\delta|x-z|}$; the separation of a primed scale-$(k-1)$ layer from the next is $D_k\defeq(r(t)-c_b)/L\ge r(t)/(2L)$ for $R\ge2c_b$, $c_b$ the large-block diameter. Two sup bounds replace every ``for $\eps_0$ small'' by a constant: for $\kappa\ge0$, $\delta>0$, $a>0$, $m\ge0$,
\[
E(\kappa,\delta)\defeq\sup_{0<t\le1}t^{-\kappa}e^{-\delta r(t)}<\infty,\qquad P(a,m)\defeq\sup_{0<t\le1}t^ap(t)^m<\infty,
\]
(with $t=e^{1-u}$, $u\ge1$: $\exp(\kappa(u-1)-\delta Ru^\varrho)$ is bounded since $\varrho>1$, and $u^{m\bar p}e^{-a(u-1)}$ is bounded), depending only on $(\kappa,\delta,R,\varrho)$ resp.\ $(a,m,b_0,\bar p)$. Three kernel facts are used throughout (sup-distance, $d(\cdot)$ a $1$-Lipschitz distance function): (K-a) if $|K(x,x')|\le ce^{-\delta(|x-x'|+d(x)+d(x'))}$ and $d\ge D$ on $Y$, then $|K|\le ce^{-2\delta D}e^{-\delta|x-x'|}$ on $Y\times Y$; (K-b) if $|K_0(x,x')|\le ce^{-\delta|x-x'|}$ for $x\in Y$, $x'\in Z$, $\dist(Y,Z)\ge D$, then $|K_0|\le ce^{-\delta D/2}e^{-\delta|x-x'|/2}$; (K-c) if $|K_1(x,z)|\le c_1e^{-\delta|x-z|}$, $|K_3(z',x')|\le c_3e^{-\delta|z'-x'|}$ and $|K_2(z,z')|\le c_2e^{-\delta(|z-z'|+d(z)+d(z'))}$, then $|(K_1K_2K_3)(x,x')|\le c_1c_2c_3S^2_{\delta/4}e^{-(\delta/4)(|x-x'|+d(x)+d(x'))}$ (with $E\defeq|x-z|+|z-z'|+|z'-x'|+d(z)+d(z')$, the triangle inequality and the Lipschitz property give $|x-x'|+d(x)+d(x')\le2E$, so the exponent $\frac{\delta}{4}(|x-x'|+d(x)+d(x'))+\frac\delta4(|x-z|+|z'-x'|)\le\delta E$; sum over $z,z'$); telescoping a product of $n$ such kernels gives the decay of (K-c) with rate $\delta/4^{n-1}$.

\begin{lemma}[the kernel types of the inverse step]\label{lem:a06d:C21}
At scale $k$ let $Y$ be a set of unit-lattice points and $D\defeq D_k$. Each of the following kernels $K(x,x')$ obeys $|K(x,x')|\le C_He^{-\delta_Hr(t)}e^{-\delta_H|x-x'|}$, with $\delta_H\defeq\min(\delta_7,\delta^S_7/4,\theta_0/4,\theta_0(\gamma_L,L)/6,\theta_W)/(16L)$ (the rates of \cite[Prop.~2.2, (2.27), (2.29)]{Ba1}, of the domain differences of $C_\Lambda$, of the Combes--Thomas bound and of the conditioning-region difference: \cref{prop:a06A:Ba1P22}, \cref{lem:a06A:domdiff}(b), \cref{lem:a06A:CT} and \cref{lem:a06A:CLdiff}) and $C_H$ explicit, all free of $k$, $\eps$, $\eps_0$ and the torus:
(a) \emph{domain change}: $K=O(\Omega_1)-O(\Omega_2)$, $\Omega_1\subset\Omega_2$ unions of big blocks, $O\in\{\Delta^{(k)}(\cdot,\tilde A),C_\Lambda(\cdot,\tilde A)\}$ at a regular background, with $\dist(Y,\Omega^{(k)c}_1)\ge D$;
(a$'$) \emph{conditioning-region change}: $K=\Delta^{(k)}_{\Lambda^{(k-1)}_5}-\Delta^{(k)}=(aL^{-2})^2Q(C-C_{\Lambda_5})Q^*$ (from \cite[(2.46)]{Ba2}, $\Delta^{(k)}_{\Lambda_5}=aL^{-2}-(aL^{-2})^2QC^{(k-1)}_{\Lambda_5}Q^*$), on $Y$ with $\dist(Y,\Omega^{(k-1)}\setminus\Lambda^{(k-1)}_5)\ge D$; this $\Lambda$-change, not covered by (a), is the only change in the photon form of \cite[(2.47), p.~568]{Ba2} not already of types (a)--(d);
(b) \emph{products}: $O_1\cdots O_n-O'_1\cdots O'_n$ ($n\le3$) with factors among $\Delta^{(k)}$, $C_\Lambda$, $a_kG_kQ_k^*$ and adjoints, each pair $O_i,O'_i$ of type (a) or equal;
(c) \emph{cut-off}: $K=\mathbf 1_YK_0\mathbf 1_Z$ with $|K_0(x,x')|\le ce^{-\delta|x-x'|}$ and $\dist(Y,Z)\ge D$;
(d) \emph{background change}: $K=O(\tilde A+\tilde B^{(k)})-O(\tilde A)$ with $|\tilde B^{(k)}|\le C_\kappa t^\kappa$ \cite[(2.90)]{Ba2}, where $e^{-\delta_Hr(t)}$ is replaced by $Ct^{\kappa+1/2}$ for every $\kappa$ (not used in Sect.~3A, which has no background change; used in Step~3(C$_2$)).
\end{lemma}

\begin{proof}
(a): \cite[(2.29)]{Ba1} (\cref{prop:a06A:Ba1P22}) gives $c_7e^{-\delta_7(|x-x'|+d(x)+d(x'))}$ with $d=\dist(\cdot,\Omega^{(k)c}_1)\ge D$ on $Y$, and the domain difference of $C_\Lambda$ the same with $\delta^S_7/4$ (\cref{lem:a06A:domdiff}(b)); then (K-a). (a$'$): the difference $C-C_{\Lambda_5}$ at scale $k-1$ decays at rate $\theta_0(\gamma_L,L)/6$ (\cref{lem:a06A:CLdiff}) in $|y-y'|+\dist(y,\Omega^{(k-1)}\setminus\Lambda^{(k-1)}_5)+\dist(y',\text{same})$, and on $\Lambda^{(k)}_5$ that distance is at least two layers; then the coarse transfer $QKQ^*$ and (K-a). (b): (K-c) and its telescoping, rate $\ge\delta/16$ for $n\le3$. (c): (K-b). In (a)--(c), $e^{-\delta D}\le e^{\delta c_b/L}e^{-\delta r(t)/L}$. (d): the analytic family of block operators in the background $\tilde A+\tau\tilde B^{(k)}$ (Step~3(C$_2$), \cref{lem:a06d:S1}) with $w=\norm qe_kC_\kappa t^\kappa$.
\end{proof}

\begin{lemma}[{\cite[(3.8), (3.10)]{Ba2}}]\label{lem:a06d:C22}
\textup{(i)} $H_k$ of \cite[p.~584]{Ba2} (``we have applied formula (2.108) together with the remark following it''): by \cite[(2.108)]{Ba2} read with the background $\tilde A^{(k)}$, the difference between $\langle\phi,\Delta^{(k)}(B^k(\Lambda^{(k-1)'}_2),\tilde A^{(k)})\phi\rangle$ on $\Lambda^{(k-1)'}_6$ and its three localised pieces is the sum of type-(a) kernels on $(\Lambda^{(k-1)'}_6\cap\Lambda^{(k)c}_3)^2$ and on $(\Lambda^{(k-1)'}_6\cap\Lambda^{(k)c}_3)\times(\Lambda^{(k)}_3\cap\Lambda^{(k)c}_4)$ ($\Omega_1=B^k(\Lambda^{(k-1)'}_2\cap\Lambda^{(k)c}_5)$), a type-(c) kernel on $(\Lambda^{(k-1)'}_6\cap\Lambda^{(k)c}_3)\times\Lambda^{(k)}_4$ (distance $\ge r-c_b$), and a type-(a) kernel on $(\Lambda^{(k)}_3)^2$ ($\Omega_1=B^k(\Lambda^{(k)}_2)\subset\Omega_2=B^k(\Lambda^{(k-1)'}_2)$); no background change occurs, since by \cite[(3.4)]{Ba2} $\tilde A^{(k)}-\tilde A^{(k+1)}=(1-\theta_{k+1})\theta_kA^{(k)}$ vanishes on $B^k(\Lambda^{(k)}_2)$. Hence \cite[(3.8)]{Ba2} holds for $h_k$ with $\delta_1=\delta_0=\delta_H$ and $O(1)=C'_H\defeq4C_H$. \textup{(ii)} $H'_k,H''_k,F'_k,F''_k$ (\cite[p.~585]{Ba2}: ``terms associated with the changes in (2.47) and the discussion following, and the term $H_k$''): by \cite[(2.47), p.~568]{Ba2} the changes are the domain splits of $\Delta^{(k)}_{\Lambda^{(k-1)}_5}(B^{k-1}(\Lambda^{(k-1)}_2),A^{(k)})$ (type (a)), the replacement of $\Delta^{(k)}_{\Lambda^{(k-1)}_5}$ by $\Delta^{(k)}$ on $\Lambda^{(k-1)'}_6$ in the third line of \cite[(2.47)]{Ba2} for the scalar and the photon form (type (a$'$)), the restriction of the third and ninth terms of \cite[(2.46)]{Ba2} to $\Lambda^{(k-1)'}_5\cap\Lambda^{(k-1)'c}_6$ (type (c)), and the Neumann conditions on $\partial B^{k-1}(\Lambda^{(k-1)}_2\cap\Lambda^{(k-1)c}_7)$ in the seventh to ninth terms, whose operators are the products $\Delta C_{\Lambda_5}\Delta$ and $\Delta C_{\Lambda_5}Q^*$ (type (b)); so $H'_k,H''_k$ obey \cite[(3.8)]{Ba2} with a fixed multiple of $C_H$. \textup{(iii)} $F'_k$ (resp.\ $F''_k$) is the sum of the type-(c) and type-(a)/(b) remainders of the linear terms and of $\sum_{x'\notin\Lambda^{(k)}_5}h'(x,x')v(x')$ with $v$ the collar fields, $|v|\le c_LM(t)$; hence $|F'_k(x)|\le C_HS_{\delta_H}e^{-\delta_Hr(t)}c_LM(t)\le C_Fe^{-\delta_Hr(t)/2}$ with $C_F\defeq c_LC_HS_{\delta_H}c_M\sup_tt^{-1}p(t)e^{-\delta_Hr(t)/2}<\infty$, and the same for $F''_k$.
\end{lemma}

\begin{proof}
Each item is the cited type with the cited separations (\cref{lem:a06d:C21}); the separations are those of the layers. Ba{\l}aban defines $H'_k,H''_k,F'_k,F''_k$ by reference (\cite[p.~568]{Ba2}: ``Terms omitted are described more precisely later''); the list of changes is his, and each is of one of the types.
\end{proof}

\begin{lemma}[{\cite[(3.8)]{Ba2} implies \cite[(3.12)]{Ba2} for every $\kappa$}]\label{lem:a06d:C23}
If $|h(x,x')|\le ce^{-\delta_1r(t)}e^{-\delta_0|x-x'|}$ on $Y$, then the operator $H$ on $\ell^2(Y)$ has $\norm H\le cS_{\delta_0}e^{-\delta_1r(t)}\le cS_{\delta_0}E(\kappa,\delta_1)t^\kappa$ for every $\kappa$.
\end{lemma}

\begin{proof}
Schur test, then the definition of $E(\kappa,\delta_1)$.
\end{proof}

\begin{lemma}[{\cite[(3.11)]{Ba2} with $\gamma_1=\frac12$, without field restriction or regularity of $\tilde A$}]\label{lem:a06d:C24}
Let $t\le\tau_1\defeq\min\{1,a_*^{1/2}/\lmu,a_*^{1/2}/m\}$. Then $G'_k\ge\frac12\lmu^2t^2$ and $G''_k\ge\frac12m^2t^2$ on $\ell^2(\Lambda^{(k)}_5)$.
\end{lemma}

\begin{proof}
$G'_k$ (resp.\ $G''_k$) is the quadratic part in $A_k|_{\Lambda_5}$ (resp.\ $\phi_k|_{\Lambda_5}$) of the first four terms of \cite[(3.9)]{Ba2}. The block-spin terms $\frac12\sum_la_{l-k}(L^{l-k})^{d-2}\sum|A_l-Q_{l-k}A_k|^2$ (resp.\ with $Q_{l-k}(\tilde A^{(k+1)})$) contribute nonnegative forms. The forms $\langle\Lambda_5A_k,\Delta^{(k)}_{\Lambda^{(k-1)}_5}\Lambda_5A_k\rangle$ and $\langle\Lambda_5\phi_k,\Delta^{(k)}_{\Lambda^{(k-1)}_5}(B^k(\Lambda^{(k-1)'}_2),\tilde A^{(k)})\Lambda_5\phi_k\rangle$ are compressions of conditional forms: the subscript records the conditional integration of \cite[(2.46)]{Ba2} over the intermediate field on $\Lambda^{(k-1)}_5$ with its collar values as separate variables, so their quadratic part is the partial elimination of that field in the two-level form (block-spin term of the step) $+\langle A',\Delta^{(k-1)}A'\rangle$, whose full elimination is $\Delta^{(k)}$ by the composition identity \cite[(2.40)--(2.41)]{Ba1}. For $H\ge0$, $m_0^2\ge0$, $a>0$ and $\norm Q\le1$, the Schur complement $\Delta\defeq a-a^2Q(H+m_0^2+aQ^*Q)^{-1}Q^*$ satisfies $\langle\psi,\Delta\psi\rangle=\min_\phi[a\norm{\psi-Q\phi}^2+\langle\phi,(H+m_0^2)\phi\rangle]\ge\min_u[a\norm{\psi-u}^2+m_0^2\norm u^2]=\frac{am_0^2}{a+m_0^2}\norm\psi^2$; a partial minimisation over a smaller set is larger, and a compression of an operator $\ge c$ is $\ge c$ on its range. The fundamental operators carry the masses $\lmu^2t^2$ (photon, $-\Delta^\eta+\lmu^2(L^k\eps)^2$) and $m^2t^2$ (scalar, $-\Delta^{\eta,N}_{A^{(k)}}+m^2(L^k\eps)^2$, \cite[p.~585]{Ba2}), and $\norm{Q_k(A)}\le1$ for every real background. So each form is $\ge a_km_0^2t^2/(a_k+m_0^2t^2)\ge\frac12m_0^2t^2$, $m_0\in\{\lmu,m\}$, since $m_0^2t^2\le a_*\le a_k$.
\end{proof}

This answers the ``corollary of Proposition~3.1'' of \cite[p.~586]{Ba2}: no restriction on the fields and no regularity of $\tilde A^{(k)}$ is needed, only $-\Delta^{\eta,N}_{A,\Omega}\ge0$ for every real background.

\begin{lemma}[{\cite[(3.15), (3.16)]{Ba2}}]\label{lem:a06d:C25}
Let a field $v_l$ on the collar $\Lambda^{(l-1)'}_{-1}\cap\Lambda^{(l)c}_5$ at scale $l\ge k+1$ obey $|v_l|\le M_v(t_l)$ (vector) or $M_s(t_l)$ (scalar), and let $v^{(k)}_l\defeq L^{-(l-k)/2}v_l(L^{-(l-k)}\cdot)$ be its image in scale-$k$ units. Then $|A^{(k)}_l|\le L^{-3(l-k)/2}M_v(t)$, $|\phi^{(k)}_l|\le L^{-3(l-k)/4}M_s(t)$ (Ba{\l}aban's $L^{-(l-k)d/2}$, $L^{-(l-k)d/4}$), and the fields $\Phi',\Phi''$ of \cite[(3.13)--(3.14)]{Ba2} satisfy $|\Phi'(x)|\le aM_v(t)/(L^{7/2}-1)$, $|\Phi''(x)|\le aM_s(t)/(L^{11/4}-1)$ on $\Lambda^{(k)}_5$, both $\le aM(t)\le ac_MP(1,1)t^{-2}$.
\end{lemma}

\begin{proof}
$t_l=L^{l-k}t$ and $p(t_l)\le p(t)$, so $L^{-(l-k)/2}c_Ap(t_l)/(\lmu t_l)\le L^{-3(l-k)/2}M_v(t)$, and likewise for the scalar. $|Q^*_{l-k}(\tilde A)\phi_l|=|\phi_l|$ (orthogonal transport) and $a_{l-k}\le a$, so $|\Phi'(x)|\le\sum_{l>k}aL^{-2(l-k)}L^{-3(l-k)/2}M_v(t)$, a geometric series; the same for $\Phi''$. Finally $M(t)\le c_MP(1,1)t^{-2}$.
\end{proof}

\begin{lemma}[the linear coefficient]\label{lem:a06d:C26}
Let $J'_k$ ($J''_k$) be the linear coefficient of $A_k|_{\Lambda_5}$ ($\phi_k|_{\Lambda_5}$) in \cite[(3.9)]{Ba2} without $F'_k$ ($F''_k$). Then $|J'_k(x)|,|J''_k(x)|\le C_JM(t)$ on $\Lambda^{(k)}_5$, hence $\norm J\le C_JM(t)|\Lambda^{(k)}_5|^{1/2}$, with $C_J\defeq a+c_{\mathrm{coll}}+c_Lc_f$: $c_{\mathrm{coll}}$ the $\ell^\infty\to\ell^\infty$ norm of the kernel of the conditional operator $\Delta^{(k)}_{\Lambda^{(k-1)}_5}$ from the collar $\Lambda^{(k-1)'}_5\cap\Lambda^{(k)c}_5$ to $\Lambda^{(k)}_5$ (its decay is that of $aL^{-2}-(aL^{-2})^2QC_{\Lambda_5}Q^*$ with $|C_{\Lambda_5}(y,y')|\le c_2e^{-\theta_0|y-y'|}$), $c_f$ that of the kernels of the third and ninth terms of \cite[(2.46)]{Ba2} (products $\Delta^{(k-1)}C^{(k-1)}_{\Lambda_5}Q^*$ of decaying kernels, $c_f\le aL^{-2}c_7c_2S_{\delta_7}S_{\theta_0}$).
\end{lemma}

\begin{proof}
$\Phi'$ by \cref{lem:a06d:C25}; the collar term by the conditional operator's decay and $|A_k|\le M_v$ on the collar; $f'_k$ is a sum of such kernels applied to scale-$(k-1)$ collar fields bounded by $c_LM(t)$.
\end{proof}

This replaces Ba{\l}aban's $O((L^k\eps)^{-2})|\Lambda^{(k)}_5|^{1/2}$ \cite[p.~587]{Ba2} by $C_Jc_Mt^{-1}p(t)|\Lambda_5|^{1/2}$.

\begin{lemma}[Gaussian comparison]\label{lem:a06d:C27}
Let $V$ be a real Euclidean space of dimension $n$, $G,H$ symmetric on $V$ with $G\ge g>0$ and $\norm H\le h\le g/2$, and $J,F\in V$. Then
\[
\int_Ve^{-\frac12\langle x,(G-H)x\rangle+\langle J+F,x\rangle}dx\le\int_Ve^{-\frac12\langle x,Gx\rangle+\langle J,x\rangle}dx\cdot\exp\Big[\frac{nh}g+\frac{h\norm J^2}{g^2}+\frac{2\norm F\norm J}g+\frac{\norm F^2}g\Big].
\]
\end{lemma}

\begin{proof}
Both integrals are Gaussian, $\int e^{-\frac12\langle x,Mx\rangle+\langle b,x\rangle}=(2\pi)^{n/2}(\det M)^{-1/2}e^{\frac12\langle b,M^{-1}b\rangle}$ for $M>0$ (here $G-H\ge g/2$). Determinant (\cite[(3.18)--(3.20)]{Ba2}): $\det(G-H)=\det G\det(I-X)$, $X\defeq G^{-1/2}HG^{-1/2}$, $\norm X\le h/g\le\frac12$; for an eigenvalue $\lambda$ of $X$, $-\log(1-\lambda)\le2|\lambda|$, so $-\frac12\log\det(I-X)\le\sum|\lambda_i|\le nh/g$. Exponent: $(G-H)^{-1}-G^{-1}=G^{-1}H(G-H)^{-1}$ has norm $\le2h/g^2$ and $\norm{(G-H)^{-1}}\le2/g$, so $\frac12\langle J+F,(G-H)^{-1}(J+F)\rangle-\frac12\langle J,G^{-1}J\rangle\le h\norm J^2/g^2+2\norm F\norm J/g+\norm F^2/g$.
\end{proof}

\begin{corollary}[{\cite[(3.9)]{Ba2} against the unperturbed integral}]\label{cor:a06d:C28}
Let $t\le\tau_1$ and assume $C'_HS_{\delta_H}e^{-\delta_HR}\le\mu_{\mathrm{lo}}^2/4$ and $\delta_HR\ge2$, $\mu_{\mathrm{lo}}^2\defeq\min(\lmu^2,m^2)$ (a condition on $R$, chosen before $\eps_0$). Then for every $\kappa>0$, \cite[(3.9)]{Ba2} is at most \cite[(3.9)]{Ba2} at $H'_k=H''_k=0$, $F'_k=F''_k=0$, times $\exp(C_{\kappa}t^\kappa|\Lambda^{(k)}_5|)$, with
\[
C_\kappa\defeq\sup_{0<t\le1}t^{-\kappa}\Big[\frac{2(d+N)h(t)}{\mu_{\mathrm{lo}}^2t^2}+\frac{4h(t)C_J^2M(t)^2}{\mu_{\mathrm{lo}}^4t^4}+\frac{4(d+N)^{1/2}C_Fe^{-\delta_Hr(t)/2}C_JM(t)}{\mu_{\mathrm{lo}}^2t^2}+\frac{2(d+N)C_F^2e^{-\delta_Hr(t)}}{\mu_{\mathrm{lo}}^2t^2}\Big],\quad h(t)\defeq C'_HS_{\delta_H}e^{-\delta_Hr(t)},
\]
finite since every term is $e^{-\delta_Hr(t)/2}$ times a power of $t^{-1}p(t)$.
\end{corollary}

\begin{proof}
$V\defeq\ell^2(\Lambda^{(k)}_5)^d\oplus\ell^2(\Lambda^{(k)}_5)^N$, $n=(d+N)|\Lambda^{(k)}_5|$; $G\defeq G'\oplus G''\ge g\defeq\frac12\mu_{\mathrm{lo}}^2t^2$ (\cref{lem:a06d:C24}); $H\defeq H'\oplus H''$ with $h=h(t)$ (\cref{lem:a06d:C22,lem:a06d:C23}), and $h\le g/2$ for all $t\le1$ by the cap ($e^{-\delta_Hr(t)}\le e^{-\delta_HR}t^2$ when $\delta_HR\ge2$); $\norm J\le C_JM(t)|\Lambda_5|^{1/2}$ (\cref{lem:a06d:C26}); $\norm F\le(d+N)^{1/2}C_Fe^{-\delta_Hr/2}|\Lambda_5|^{1/2}$ (\cref{lem:a06d:C22}(iii)). \cref{lem:a06d:C27}, with the scalar background treated as an independent external field (\cite[p.~568]{Ba2}), gives the bound.
\end{proof}

\begin{lemma}[{the collar forms, \cite[p.~584]{Ba2}}]\label{lem:a06d:C29}
On the collar $\Lambda^{(k-1)'}_5\cap\Lambda^{(k)c}_5$, where $|A_k|\le M_v$, $|\phi_k|\le M_s$, every form $\frac12\langle v,Kv\rangle$ with $K$ of the types of \cref{lem:a06d:C21} obeys $|\frac12\langle v,Kv\rangle|\le\frac12C'_HS_{\delta_H}e^{-\delta_Hr(t)}M(t)^2|\text{collar}|\le C'_\kappa t^\kappa|\Lambda^{(k-1)'}_5\cap\Lambda^{(k)c}_5|$ with $C'_\kappa\defeq\frac12C'_HS_{\delta_H}c_L^2c_M^2\sup_tt^{-\kappa-2}p(t)^2e^{-\delta_Hr(t)}<\infty$ (Schur test).
\end{lemma}

\begin{proposition}[{the inductive step \cite[(3.5)]{Ba2}$_k\Rightarrow$\cite[(3.5)]{Ba2}$_{k-1}$}]\label{prop:a06d:C211}
Assume \cite[(3.5)]{Ba2} at level $k$. Then it holds at level $k-1$ with the additional factor $\exp(C_{\mathrm{LF},\kappa}t_k^\kappa|T^{(k)}_1|)$, $C_{\mathrm{LF},\kappa}\defeq C_\kappa+3C'_\kappa$: the $T$-transformations are integral operators with nonnegative kernels, so inequalities inside the nested brackets are preserved; the transformations not depending on $A_k,\phi_k|_{\Lambda^{(k)}_5}$ are moved out (Fubini--Tonelli); the forms of \cite[(3.5)]{Ba2} at level $k$ are merged by \cite[(2.108)]{Ba2} into one form plus $H_k$ (\cref{lem:a06d:C22}(i)), the forms evaluated on the collar cost \cref{lem:a06d:C29}, the remaining integral is \cite[(3.9)]{Ba2} (\cref{lem:a06d:C22}(ii)) and \cref{cor:a06d:C28} compares it with the $H=F=0$ integral; the latter is, by the definitions \cite[(2.45)--(2.50)]{Ba2} read backwards, the level-$k$ factor of $\varrho''^{(k)}$ \cite[(2.46)]{Ba2}, and \cite[(2.48)]{Ba2} replaces $\varrho''$ by $\varrho$ (for $k<K$ its interior part is of $H'$-type and sits inside \cref{cor:a06d:C28}, its collar part is a \cref{lem:a06d:C29} form; at $k=K$, where \cite[p.~583]{Ba2} restricts $B,\psi$ everywhere, the whole of \cite[(2.48)]{Ba2} is a sum of \cref{lem:a06d:C29} forms, which is the base case \cite[(3.1)]{Ba2}). The sum over levels of the added factors is $\sum_lC_{\mathrm{LF},\kappa}t_l^\kappa|T^{(l)}_1|=\sum_lC_{\mathrm{LF},\kappa}t_l^{\kappa-d}|T_\eps|$, Ba{\l}aban's $\exp(\sum_lO(1)(L^l\eps)^{\kappa_0}|T_\eps|)$ with $\kappa_0=\kappa-d$ and $O(1)=C_{\mathrm{LF},\kappa}$ independent of $l$; the induction is finite ($K$ steps).
\end{proposition}

Every constant above is explicit and free of $k$, $\eps$, $\eps_0$ and the torus; the caps used are $t\le\tau_1$ and the condition on $R$ of \cref{cor:a06d:C28}, neither involving $\eps_0$.
